In a controlled bimodal reach task, an MSE head fails when starts are nearly identical (success \rhval{sweep_jitter/mse-mlp-jit-0.0/bimodal50/success/mean:pct1}\% at zero jitter) but recovers as start jitter grows, and on skewed demonstrations it succeeds by collapsing to the majority mode. A two-component Gaussian mixture and a 50-step DDPM head both solve the task and reproduce the demonstrated mode share; we found no advantage of DDPM over the mixture, and the DDPM head was weaker with one denoising step and with long open-loop chunks. The result is narrow: one toy environment, small MLPs, reimplemented untuned baselines. A fair test of diffusion heads needs tasks where a mixture is not at ceiling, such as more modes, higher-dimensional or sequential multimodality.
